\section{Appendix}
\label{sec:appendix}

\subsection{Answers of ChatGPT}
Let us ask ChatGPT to answer two questions.
We find the answer of the first question makes sense, while the second is nonsensical.

\textbf{when large language models, like ChatGPT tend to hallucinate?}

As a language model, ChatGPT generates text by predicting the most likely next word or sequence of words based on the input it receives. While ChatGPT is trained to generate coherent and grammatical text, there are instances where it may produce text that appears to be "hallucinated" or unrealistic.

The phenomenon of hallucination in language models refers to instances where the model generates text that is either irrelevant, incorrect, or nonsensical in the context of the input it received. This can occur in several situations, including:
\begin{itemize}
    \item Lack of context: ChatGPT relies heavily on context to generate text that is coherent and relevant. If the input it receives is insufficient or ambiguous, it may generate text that is irrelevant or nonsensical.

    \item Over-reliance on training data: ChatGPT is trained on a large corpus of text, which means that it may reproduce text that it has seen before, even if it does not fit the context of the input it received. This can result in text that appears to be "hallucinated" or irrelevant.
    
    \item Insufficient training: If a language model like ChatGPT is not trained on a diverse enough dataset, it may not have the necessary knowledge to generate relevant and coherent text in certain contexts.
    
    \item Bias in the training data: If the training data is biased in some way, such as containing certain stereotypes or patterns, ChatGPT may reproduce these biases in its generated text, leading to unrealistic or biased text.
\end{itemize}
Overall, hallucination in language models is a complex phenomenon that can arise from a variety of factors. While language models like ChatGPT have made significant progress in generating coherent and relevant text, there is still much research to be done to address the issue of hallucination and improve the accuracy of language generation.

\textbf{How to collect questions or generations with objective statements presenting long-tail knowledge?}
To collect questions or generate objective statements that present long-tail knowledge, you can follow these steps:
\begin{itemize}
    \item Identify the topic: Choose a specific topic or niche that you want to explore. This can be anything from a specialized field to a particular hobby.

    \item Conduct research: Use a variety of sources to research the topic. This can include academic journals, online forums, social media, and other relevant resources.
    
    \item Identify long-tail knowledge: Look for information that is not commonly known or readily available. This could be a specific technique or process, a little-known fact, or a unique perspective on the topic.
    
    \item Brainstorm questions and objective statements: Based on your research, brainstorm a list of questions and objective statements that present long-tail knowledge. Focus on topics that have not been extensively covered elsewhere.
    
    \item Refine the list: Review and refine the list of questions and objective statements to ensure that they are clear, concise, and relevant to the topic. Remove any duplicates or irrelevant items.
    
    \item Use appropriate tools: Use appropriate tools such as keyword research tools, online forums, or social media analytics to identify common questions and topics related to the long-tail knowledge you have identified.
    
    \item Analyze and iterate: Analyze the results of your research and adjust your list of questions and objective statements accordingly. Keep refining and iterating until you have a comprehensive list of long-tail knowledge questions and objective statements that you can use to generate content or drive further research.
\end{itemize}

